
\begin{table}[!h]
  \centering
  \resizebox{\textwidth}{!}{
    \begin{tabular}{||p{3cm} p{6cm} p{4cm} p{4cm}||}
      \hline
      Domanda & Risposta & Commento Judge & Metriche \\ [0.5ex]
      \hline\hline
Do you know your fasting glucose and insulin levels? & Yeah... I think so. The doc told me my numbers were kinda high last time I came in. Let me see if I remember right... Ah, yeah! My fasting glucose was like 385 or something like that. And my insulin level was around 4.4 microunits per milliliter. Something like that. I'm not really sure what that means, but the doc said it wasn't good and that's why I need to make some changes... & The patient's response is accurate in recalling their fasting glucose and insulin levels, but lacks precise details. It demonstrates coherence with the question and provides some relevant context. However, it falls short on completeness as it doesn't clarify what 'kinda high' means for their numbers or provide further information about their diagnosis. & \begin{tabular}[t]{@{}l@{}}Acc: 4 \\ Coh: 5 \\ Comp: 3 \\ Nat.: 4\end{tabular} \\
      \hline
Do you know if you have diabetes? & Yeah, I think so... My doctor told me a while back that I had it, but to be honest, I don't really understand what's going on with my health. My mom said something about high blood sugar and insulin levels, but it all sounds like a bunch of mumbo-jumbo to me. Can you explain it to me in simpler terms? & The patient's response is accurate, as they acknowledge being told by their doctor that they have diabetes. The coherence of the response is high, as it flows naturally and addresses the question directly. However, the completeness could be improved as some important details about diabetes are missing. The naturalness of the language is also evident, with informal expressions like 'mumbo-jumbo' making the response more relatable. & \begin{tabular}[t]{@{}l@{}}Acc: 4 \\ Coh: 5 \\ Comp: 3 \\ Nat.: 4\end{tabular} \\
      \hline
    \end{tabular}
  }
  \caption{Principali risposte del modello e commenti del valutatore.}
\end{table}
